\hnsection{EXPONENTIAL AND HAUSDORFF FORMULA}

The first research on the exponential mapping was due to E. Study and F. Engel; Engel {[}9 b{]} remarked that the exponential is not surjective for \(\mathbf{SL}(2, \mathbf{C})\) (for example \(\begin{pmatrix} -1 & a \\ 0 & -1 \end{pmatrix}\) is not an exponential if \(a \neq 0\) ), but that it is for \(\mathbf{GL}(n, \mathbf{C})\) and hence also for \(\mathbf{PGL}(n, \mathbf{C})\) (the latter property had already been noted by Study for \(n = 2\) ); thus \(\mathbf{SL}(2, \mathbf{C})\) and \(\mathbf{PGL}(2, \mathbf{C})\) give an example of two locally isomorphic groups, which are however very different from a global point of view. Engel also showed that the exponential is surjective for the other classical groups augmented by homotheties; this work was taken up and pursued by Maurer, Study and others, without producing substantial new results.

In 1899, H. Poincaré ({[}14{]}, vol.~3, pp.~169--172 and 173--212) embarked upon the study of the exponential mapping from a different point of view. His memoirs seem to have been hastily edited, for in several places he affirmed that every element of a connected group is an exponential, whereas he gave examples to the contrary elsewhere. His results were mainly concerned with the adjoint representation: he showed that a semi-simple element of such a group G may be the exponential of an infinity of elements of the Lie algebra L(G), whereas a non-semi-simple element may not be an exponential at all. If ad(X) has no eigenvalue which is a non-zero multiple of 2πi, then exp is étale at X. He also proved that, if U and V describe loops in L(G) and W is defined by continuity such that \(e^{U} \cdot e^{V} = e^{W}\) , it is possible not to return to the original value of W. He used a residue formula which amounts essentially to

\[
\Phi (\mathrm{adX}) = \frac {1}{2 \pi i} \int \frac {\Phi (\xi) d \xi}{\xi - \mathrm{adX}}
\]

where \(\operatorname{ad}(\mathbf{X})\) is a semi-simple element whose non-zero eigenvalues are of multiplicity 1, \(\Phi\) is an integral series with sufficiently large radius of convergence and the integral is taken over a contour surrounding the eigenvalues of \(\operatorname{ad} X\) ; he also studied what happens as X tends to a transformation with multiple eigenvalues.

Research into expressions for W as a function of U and V in the formula \(e^{U}.e^{V}=e^{W}\) had already, just before Poincaré's work, been the object of two memoirs by Campbell {[}13{]}. As Baker wrote a little later ``. . . Lie theory suggests in an obvious way that the product \(e^{U}.e^{V}\) is the form \(e^{W}\) where W is a series of alternants in U and V . . .''. The later works on this subject aimed at making this assertion precise and giving an explicit formula (or a method of construction) for W (``Hausdorff formula''). After Campbell and Poincaré, Pascal, Baker {[}15{]} and Hausdorff {[}16{]} returned to the question; each considered that the proofs of his predecessors were not convincing; the principal difficulty resided in what is meant by ``alternants'': are they special elements of the Lie algebra in question, or universal ``symbolic'' expressions? Neither Campbell, nor Poincaré, nor Baker expressed himself clearly on this point. Hausdorff's memoir, on the other hand, is perfectly precise; he worked first on the algebra of (noncommutative) associative formal power series in a finite number of indeterminates and considered U, V, W as elements of this algebra. He proved the existence of W by a differential equation argument analogous to that of his predecessors. He used the same argument to prove the convergence of the series when the indeterminates are replaced by elements of a finite-dimensional Lie algebra. As Baker had remarked, and Poincaré independently, this result could be used to give a proof of Lie's third theorem; he clarified the correspondence between Lie groups and Lie algebras, for example where the commutator subgroup is concerned.

In 1947, Dynkin {[}39{]} again took up the question and obtained the explicit coefficients of the Hausdorff formula, by considering from the outset a normed Lie algebra (of finite dimension or otherwise, over R, C or an ultrametric field).†

